\chapter{2010年安徽卷（理）}

\section{单选题}

\begin{problem}
\(\i\) 是虚数单位，\( \frac{\i}{\sqrt{3} + 3\i} = \)（\quad）

\choices
    {\(\frac{1}{4} - \frac{\sqrt{3}}{12}\i\)}
    {\(\frac{1}{4} + \frac{\sqrt{3}}{12}\i\)}
    {\(\frac{1}{2} + \frac{\sqrt{3}}{6}\i\)}
    {\(\frac{1}{2} - \frac{\sqrt{3}}{6}\i\)}
\end{problem}

\begin{answer}
B
\end{answer}

\begin{solution}
将分子、分母同乘以分母的共轭复数，得
\[
\frac{\i}{\sqrt{3}+3\i}
=\frac{\i(\sqrt{3}-3\i)}{(\sqrt{3}+3\i)(\sqrt{3}-3\i)}
=\frac{3+\sqrt{3}\i}{12}
=\frac{1}{4}+\frac{\sqrt{3}}{12}\i
\]
因此选 B.
\end{solution}

\begin{problem}
若集合 \(A = \left\{x \mid \log_{\frac{1}{2}} x \geqslant \frac{1}{2}\right\}\)，则 \(\complement_{\R} A =\)（\quad）

\choices
    {\((-\infty, 0] \cup \left(\frac{\sqrt{2}}{2}, +\infty\right)\)}
    {\(\left(\frac{\sqrt{2}}{2}, +\infty\right)\)}
    {\((-\infty, 0] \cup \left[\frac{\sqrt{2}}{2}, +\infty\right)\)}
    {\(\left[\frac{\sqrt{2}}{2}, +\infty\right)\)}
\end{problem}

\begin{answer}
A
\end{answer}

\begin{solution}
对数的定义域要求 \(x>0\). 因为底数 \(\frac{1}{2}\) 在 \(0\) 与 \(1\) 之间，故不等号方向改变，
\[
\log_{\frac{1}{2}}x\geqslant\frac{1}{2}
\Longleftrightarrow
0<x\leqslant\left(\frac{1}{2}\right)^{\frac{1}{2}}=\frac{\sqrt{2}}{2}
\]
所以 \(A=(0,\frac{\sqrt{2}}{2}]\)，从而
\[
\complement_{\R}A=(-\infty,0]\cup\left(\frac{\sqrt{2}}{2},+\infty\right)
\]
因此选 A.
\end{solution}

\begin{problem}
设向量 \(\bs{a} = (1,0)\)，\(\bs{b} = \left(\frac{1}{2}, \frac{1}{2}\right)\)，则下列结论中正确的是（\quad）

\choices
    {\(|\bs{a}| = |\bs{b}|\)}
    {\(\bs{a} \cdot \bs{b} = \frac{\sqrt{2}}{2}\)}
    {\(\bs{a} - \bs{b}\) 与 \(\bs{b}\) 垂直}
    {\(\bs{a} \parallel \bs{b}\)}
\end{problem}

\begin{answer}
C
\end{answer}

\begin{solution}
由题设
\(|\bs{a}|=\sqrt{1^2+0^2}=1\)，\(|\bs{b}|=\sqrt{\left(\frac{1}{2}\right)^2+\left(\frac{1}{2}\right)^2}=\frac{\sqrt{2}}{2}\)
故选项 A 错误. 又
\[
\bs{a}\cdot\bs{b}=1\cdot\frac{1}{2}+0\cdot\frac{1}{2}=\frac{1}{2}\neq\frac{\sqrt{2}}{2}
\]
故选项 B 错误. 此外
\(\bs{a}-\bs{b}=\left(\frac{1}{2},-\frac{1}{2}\right)\)，\((\bs{a}-\bs{b})\cdot\bs{b}=\frac{1}{4}-\frac{1}{4}=0\)
所以 \(\bs{a}-\bs{b}\) 与 \(\bs{b}\) 垂直. 两向量的对应坐标不成比例，故不平行，选 C.
\end{solution}

\begin{problem}
若 \( f(x) \) 是 \( \R \) 上周期为 5 的奇函数，且满足 \(f(1) = 1\)，\(f(2) = 2\)，则 \( f(3) - f(4) = \)（\quad）

\choices
    {\(-1\)}
    {\(1\)}
    {\(-2\)}
    {\(2\)}
\end{problem}

\begin{answer}
A
\end{answer}

\begin{solution}
由周期性，\(f(3)=f(-2)\)，\(f(4)=f(-1)\). 又 \(f\) 为奇函数，所以
\(f(3)=f(-2)=-f(2)=-2\)，\(f(4)=f(-1)=-f(1)=-1\).
因此 \(f(3)-f(4)=-2-(-1)=-1\)，故选 A.
\end{solution}

\begin{problem}
双曲线方程为 \(x^{2} - 2y^{2} = 1\)，则它的右焦点坐标为（\quad）

\choices
    {\(\left(\frac{\sqrt{2}}{2}, 0\right)\)}
    {\(\left(\frac{\sqrt{5}}{2}, 0\right)\)}
    {\(\left(\frac{\sqrt{6}}{2}, 0\right)\)}
    {\((\sqrt{3},0)\)}
\end{problem}

\begin{answer}
C
\end{answer}

\begin{solution}
将双曲线方程化为标准形式：
\[
\frac{x^2}{1^2}-\frac{y^2}{\left(\frac{\sqrt{2}}{2}\right)^2}=1
\]
因此 \(a^2=1\)，\(b^2=\frac{1}{2}\). 双曲线的焦半距满足
\(c^2=a^2+b^2=1+\frac{1}{2}=\frac{3}{2}\)，\(c=\frac{\sqrt{6}}{2}\).
右焦点为 \(\left(\frac{\sqrt{6}}{2},0\right)\)，故选 C.
\end{solution}

\begin{problem}
设 \( abc > 0 \)，二次函数 \( f(x) = ax^2 + bx + c \) 的图象可能是（\quad）

\choices
    {\choicebitmap[width=0.15\paperwidth]{img/2010/anhui_science/q6_optA.png}}
    {\choicebitmap[width=0.15\paperwidth]{img/2010/anhui_science/q6_optB.png}}
    {\choicebitmap[width=0.15\paperwidth]{img/2010/anhui_science/q6_optC.png}}
    {\choicebitmap[width=0.15\paperwidth]{img/2010/anhui_science/q6_optD.png}}
\end{problem}

\begin{answer}
D
\end{answer}

\begin{solution}
由 \(abc>0\)，且 \(a\)，\(b\)，\(c\) 均非零，需结合图象开口方向、\(y\) 轴截距和对称轴位置判断.

图象 A 中 \(a<0\)，\(c<0\)，若 \(abc>0\)，则必须有 \(b>0\)，此时对称轴
\(x=-\frac{b}{2a}>0\)，与图象 A 的对称轴在 \(y\) 轴左侧矛盾. 图象 B 中 \(a<0\)，\(c>0\)，故必须 \(b<0\)，从而 \(x=-\frac{b}{2a}<0\)，与图象 B 的对称轴在 \(y\) 轴右侧矛盾. 图象 C 中 \(a>0\)，\(c<0\)，故必须 \(b<0\)，从而对称轴应在 \(y\) 轴右侧，与图象 C 矛盾.

图象 D 中 \(a>0\)，\(c<0\)，\(b<0\)，所以 \(abc>0\)，且其开口、截距和对称轴位置均相符. 因此选 D.
\end{solution}

\begin{problem}
设曲线 \(C\) 的参数方程为 \(\begin{cases}
x = 2 + 3\cos \theta,\\
y = -1 + 3\sin \theta,
\end{cases}\) （\(\theta\) 为参数），直线 \(l\) 的方程为 \(x - 3y + 2 = 0\)，则曲线 \(C\) 上到直线 \(l\) 的距离为 \(\frac{7\sqrt{10}}{10}\) 的点的个数为（\quad）

\choices
    {\(1\)}
    {\(2\)}
    {\(3\)}
    {\(4\)}
\end{problem}

\begin{answer}
B
\end{answer}

\begin{solution}
曲线 \(C\) 上点的坐标满足
\(x=2+3\cos\theta\)，\(y=-1+3\sin\theta\).
代入直线 \(l\) 的左端，得
\[
x-3y+2=7+3\cos\theta-9\sin\theta
\]
点到直线 \(l\) 的距离为 \(\frac{7\sqrt{10}}{10}=\frac{7}{\sqrt{10}}\)，所以
\[
\left|7+3\cos\theta-9\sin\theta\right|=7
\]
于是有
\[
3\cos\theta-9\sin\theta=0
\quad\text{或}\quad
3\cos\theta-9\sin\theta=-14
\]
而
\[
\left|3\cos\theta-9\sin\theta\right|\leqslant\sqrt{3^2+(-9)^2}=3\sqrt{10}<14
\]
第二种情况无解. 第一种情况化为 \(\tan\theta=\frac{1}{3}\)，在一个周期内有两个不同的参数值，且对应两个不同的圆上点. 所以点的个数为 \(2\)，选 B.
\end{solution}

\begin{problem}
一个几何体的三视图如图，该几何体的表面积为（\quad）

\bitmapfigure[max width=0.74\linewidth]{img/2010/anhui_science/q8_fig1.png}

\choices
    {\(280\)}
    {\(292\)}
    {\(360\)}
    {\(372\)}
\end{problem}

\begin{answer}
C
\end{answer}

\begin{solution}
由三视图可知，下面的长方体长为 \(8=6+1+1\)，宽为 \(10=6+2+2\)，高为 \(2\)；上面的长方体长为 \(6\)，宽为 \(2\)，高为 \(8\). 两个长方体的接触面为 \(6\times2\).

接触面面积为 \(6\times2=12\). 下面长方体可见的表面积为
\[
2(8\times10+8\times2+10\times2)-6\times2=220
\]
上面长方体露出的表面积包括四个侧面和上底面，等于
\[
2(6+2)\times8+6\times2=140
\]
所以几何体的表面积为 \(220+140=360\)，故选 C.
\end{solution}

\begin{problem}
动点 \(A(x, y)\) 在圆 \(x^2 + y^2 = 1\) 上绕坐标原点沿逆时针方向匀速旋转，\(12\text{ 秒}\)旋转一周. 已知时间 \(t = 0\) 时，点 \(A\) 的坐标是 \(\left(\frac{1}{2}, \frac{\sqrt{3}}{2}\right)\)，则当 \(0 \leqslant t \leqslant 12\) 时，动点 \(A\) 的纵坐标 \(y\) 关于 \(t\)（单位：秒）的函数的单调递增区间是（\quad）

\choices
    {\([0,1]\)}
    {\([1,7]\)}
    {\([7,12]\)}
    {\([0,1]\) 和 \([7,12]\)}
\end{problem}

\begin{answer}
D
\end{answer}

\begin{solution}
初始时点 \(A\) 的极角为 \(\frac{\pi}{3}\)，每秒转过的角为
\[
\omega=\frac{2\pi}{12}=\frac{\pi}{6}
\]
所以
\(y=\sin\left(\frac{\pi}{3}+\frac{\pi}{6}t\right)\)，\(0\leqslant t\leqslant12\).
其导数为
\[
y'=\frac{\pi}{6}\cos\left(\frac{\pi}{3}+\frac{\pi}{6}t\right)
\]
因此 \(y\) 单调递增当且仅当余弦值为正. 在 \(0\leqslant t\leqslant12\) 内，这给出
\[
0\leqslant t\leqslant1\quad\text{或}\quad 7\leqslant t\leqslant12
\]
故选 D.
\end{solution}

\begin{problem}
设 \(\{a_{n}\}\) 是任意等比数列，它的前 \(n\) 项和，前 \(2n\) 项和与前 \(3n\) 项和分别为 \(X\)，\(Y\)，\(Z\)，则下列等式中恒成立的是（\quad）

\choices
    {\(X + Z = 2Y\)}
    {\(Y(Y - X) = Z(Z - X)\)}
    {\(Y^{2} = XZ\)}
    {\(Y(Y - X) = X(Z - X)\)}
\end{problem}

\begin{answer}
D
\end{answer}

\begin{solution}
设等比数列的公比为 \(q\)，令第一组 \(n\) 项的和为 \(X\). 第二组 \(n\) 项是第一组各项的 \(q^n\) 倍，第三组 \(n\) 项是第一组各项的 \(q^{2n}\) 倍，因此
\(Y=X(1+q^n)\)，\(Z=X(1+q^n+q^{2n})\).
于是
\(Y-X=Xq^n\)，\(Z-X=Xq^n(1+q^n)\).
从而
\[
Y(Y-X)=X(1+q^n)\cdot Xq^n
=X\cdot Xq^n(1+q^n)=X(Z-X)
\]
所以恒成立的等式为 D.
\end{solution}

\section{填空题}

\begin{problem}
命题“对任何 \(x \in \R\)，\(|x - 2| + |x - 4| > 3\)”的否定是\fillinblank{}.
\end{problem}

\begin{answer}
存在 \(x \in \R\)，使得 \(|x - 2| + |x - 4| \leqslant 3\)
\end{answer}

\begin{solution}
原命题是全称命题，对任意 \(x\in\R\) 有 \(|x-2|+|x-4|>3\). 否定全称量词要改为存在量词，同时将 \(>3\) 否定为 \(\leqslant3\). 因此否定为：存在 \(x\in\R\)，使得
\[
|x-2|+|x-4|\leqslant3
\]
\end{solution}

\begin{problem}
\(\left(\frac{x}{\sqrt{y}} - \frac{y}{\sqrt{x}}\right)^{6}\) 的展开式中，\(x^{3}\) 的系数等于\fillinblank{}.
\end{problem}

\begin{answer}
15
\end{answer}

\begin{solution}
根式在分母中，所以 \(x>0\)，\(y>0\). 按二项式定理展开，通项为
\[
\mathrm C_6^k\left(\frac{x}{\sqrt y}\right)^{6-k}
\left(-\frac{y}{\sqrt x}\right)^k
=(-1)^k\mathrm C_6^k x^{6-\frac{3k}{2}}y^{-3+\frac{3k}{2}}
\]
要使 \(x\) 的指数为 \(3\)，必须有
\[
6-\frac{3k}{2}=3\Longrightarrow k=2
\]
此时 \(y\) 的指数为 \(0\)，所以 \(x^3\) 的系数为
\[
(-1)^2\mathrm C_6^2=15
\]
\end{solution}

\begin{problem}
设 \(x\)，\(y\) 满足约束条件 \(\begin{cases}
2x - y + 2 \geqslant 0,\\
8x - y - 4 \leqslant 0,\\
x \geqslant 0,\\
y \geqslant 0
\end{cases}\)，
若目标函数 \(z = abx + y\)（\(a > 0\)，\(b > 0\)）的最大值为 8，则 \(a + b\) 的最小值为\fillinblank{}.
\end{problem}

\begin{answer}
4
\end{answer}

\begin{solution}
约束条件等价于
\(y\leqslant2x+2\)，\(y\geqslant8x-4\)，\(x\geqslant0\)，\(y\geqslant0\).
可行域的顶点为 \((0,0)\)、\((\frac{1}{2},0)\)、\((1,4)\)、\((0,2)\). 令 \(p=ab>0\)，在这些顶点处目标函数 \(z=px+y\) 的值分别为
\(0\)，\(\frac{p}{2}\)，\(p+4\)，\(2\).
因为 \(p>0\)，最大值为 \(p+4\). 由最大值为 \(8\)，得 \(p=4\)，即 \(ab=4\). 再由算术—几何平均不等式，
\[
a+b\geqslant2\sqrt{ab}=4
\]
且当 \(a=b=2\) 时取等号. 所以 \(a+b\) 的最小值为 \(4\).
\end{solution}

\begin{problem}
如图所示，程序框图（算法流程图）的输出值 \(x =\fillinblank{}\).

\texfigure[max width=0.28\linewidth]{img_repaint/2010/anhui_science/q14_fig1.tex}
\end{problem}

\begin{answer}
12
\end{answer}

\begin{solution}
程序开始时 \(x=1\). 逐次沿流程图执行：\(1\) 是奇数，变为 \(2\)；\(2\) 不是奇数，变为 \(4\)，且 \(4\not>8\)，再加 \(1\) 变为 \(5\)；\(5\) 是奇数，变为 \(6\)；\(6\) 不是奇数，变为 \(8\)，且 \(8\not>8\)，再加 \(1\) 变为 \(9\)；\(9\) 是奇数，变为 \(10\)；\(10\) 不是奇数，变为 \(12\)，此时 \(12>8\)，输出 \(12\).
\end{solution}

\begin{problem}
甲罐中有5个红球，2个白球和3个黑球，乙罐中有4个红球，3个白球和3个黑球，先从甲罐中随机取出一球放入乙罐，分别以 \(A_1\)，\(A_2\) 和 \(A_3\) 表示由甲罐取出的球是红球，白球和黑球的事件，再从乙罐中随机取出一球，以 \(B\) 表示由乙罐取出的球是红球的事件，则下列命题中正确的是\fillinblank{}.（写出所有正确结论的编号）

\begin{circlelist}
\item \(P(B)=\frac{2}{5}\)；
\item \(P(B \mid A_1)=\frac{5}{11}\)；
\item 事件 \(B\) 与事件 \(A_1\) 相互独立；
\item \(A_1\)，\(A_2\)，\(A_3\) 是两两互斥的事件；
\item \(P(B)\) 的值不能确定，因为它与 \(A_1\)，\(A_2\)，\(A_3\) 中究竟哪一个发生有关.
\end{circlelist}
\end{problem}

\begin{answer}
\(\circled{2}\circled{4}\)
\end{answer}

\begin{solution}
三个事件 \(A_1\)，\(A_2\)，\(A_3\) 分别表示从甲罐取出红球、白球、黑球，因此
\(P(A_1)=\frac{5}{10}=\frac{1}{2}\)，\(P(A_2)=\frac{2}{10}=\frac{1}{5}\)，\(P(A_3)=\frac{3}{10}\).
若发生 \(A_1\)，乙罐中有 \(5\) 个红球，共 \(11\) 个球，所以
\[
P(B\mid A_1)=\frac{5}{11}
\]
若发生 \(A_2\) 或 \(A_3\)，乙罐中均有 \(4\) 个红球，共 \(11\) 个球，所以
\[
P(B\mid A_2)=P(B\mid A_3)=\frac{4}{11}
\]
由全概率公式，
\[
P(B)=\frac{1}{2}\cdot\frac{5}{11}+\frac{1}{5}\cdot\frac{4}{11}
+\frac{3}{10}\cdot\frac{4}{11}=\frac{9}{22}
\]
故结论 \circled{1} 错误，结论 \circled{2} 正确. 又
\(P(B\cap A_1)=P(A_1)P(B\mid A_1)=\frac{5}{22}\)，\(P(B)P(A_1)=\frac{9}{44}\)
二者不相等，所以 \(B\) 与 \(A_1\) 不独立，结论 \circled{3} 错误. \(A_1\)，\(A_2\)，\(A_3\) 分别对应一次取球的三种互不相容结果，任意两个不能同时发生，故结论 \circled{4} 正确. \(P(B)\) 已由全概率公式唯一确定，故结论 \circled{5} 错误. 所有正确结论的编号为 \(2,4\).
\end{solution}

\section{解答题}

\begin{problem}
设 \(\triangle ABC\) 是锐角三角形，\(a\)，\(b\)，\(c\) 分别是内角 \(A\)，\(B\)，\(C\) 所对边长，并且 \(\sin^{2} A = \sin \left(\frac{\pi}{3} + B\right) \sin \left(\frac{\pi}{3} - B\right) + \sin^{2} B\).

\begin{enumerate}
\item 求角 \(A\) 的值；
\item 若 \(\overrightarrow{AB} \cdot \overrightarrow{AC} = 12\)，\(a = 2 \sqrt{7}\)，求 \(b\)，\(c\)（其中 \(b < c\)）.
\end{enumerate}
\end{problem}

\begin{solution}
\begin{enumerate}
\item 利用公式 \(\sin(u+v)\sin(u-v)=\sin^2u-\sin^2v\)，原等式右端为
\[
\sin^2\frac{\pi}{3}-\sin^2B+\sin^2B=\frac{3}{4}
\]
所以 \(\sin^2A=\frac{3}{4}\). 因为三角形为锐角三角形，\(0<A<\frac{\pi}{2}\)，故
\[
A=\frac{\pi}{3}
\]
\item 向量 \(\overrightarrow{AB}\) 与 \(\overrightarrow{AC}\) 的夹角为 \(A\)，且 \(|\overrightarrow{AB}|=c\)，\(|\overrightarrow{AC}|=b\)，所以
\[
12=\overrightarrow{AB}\cdot\overrightarrow{AC}=bc\cos A=\frac{bc}{2}
\]
从而 \(bc=24\). 由余弦定理，
\[
a^2=b^2+c^2-2bc\cos A=b^2+c^2-24
\]
代入 \(a=2\sqrt7\)，得 \(b^2+c^2=52\). 因此
\[
(b+c)^2=b^2+c^2+2bc=52+48=100
\]
由于 \(b\)，\(c>0\)，有 \(b+c=10\). 所以 \(b\)，\(c\) 是方程
\[
u^2-10u+24=0
\]
的两个根，即 \(u=4\) 或 \(u=6\). 结合 \(b<c\)，得 \(b=4\)，\(\ c=6\).
\end{enumerate}
\end{solution}

\begin{problem}
设 \(a\) 为实数，函数 \(f(x) = e^{x} - 2x + 2a\)，\(x \in \R\).

\begin{enumerate}
\item 求 \(f(x)\) 的单调区间与极值；
\item 证明：当 \(a > \ln 2 - 1\) 且 \(x > 0\) 时，\(e^{x} > x^{2} - 2ax + 1\).
\end{enumerate}
\end{problem}

\begin{solution}
\begin{enumerate}
\item
\[
f'(x)=e^x-2
\]
当 \(x<\ln2\) 时，\(f'(x)<0\)；当 \(x>\ln2\) 时，\(f'(x)>0\). 所以 \(f(x)\) 在 \((-\infty,\ln2)\) 上单调递减，在 \((\ln2,+\infty)\) 上单调递增，并在 \(x=\ln2\) 处取得极小值
\[
f(\ln2)=2-2\ln2+2a=2(1-\ln2+a)
\]
又因为 \(x\to-\infty\) 时 \(-2x\to+\infty\)，\(x\to+\infty\) 时 \(e^x\to+\infty\)，所以 \(f(x)\) 无极大值.
\item 令
\[
g(x)=e^x-x^2+2ax-1
\]
则
\[
g'(x)=e^x-2x+2a=f(x)
\]
由第（1）问，\(f(x)\) 的最小值为 \(2(1-\ln2+a)>0\)，所以 \(g'(x)>0\) 对一切 \(x\in\R\) 成立，\(g\) 在 \(\R\) 上单调递增. 当 \(x>0\) 时，
\[
g(x)>g(0)=1-0+0-1=0
\]
于是 \(e^x-x^2+2ax-1>0\)，即
\[
e^x>x^2-2ax+1
\]
\end{enumerate}
\end{solution}

\begin{problem}
如图，在多面体 \(ABCDEF\) 中，四边形 \(ABCD\) 是正方形，\(EF \parallel AB\)，\(EF \perp FB\)，\(AB = 2EF\)，\(\angle BFC = 90^\circ\)，\(BF = FC\)，\(H\) 为 \(BC\) 的中点.

\begin{enumerate}
\item 求证： \(FH \parallel\) 平面 \(EDB\)；
\item 求证： \(AC \perp\) 平面 \(EDB\)；
\item 求二面角 \(B - DE - C\) 的大小.
\end{enumerate}

\FigureLayoutDeclare{img/2010/anhui_science/q18_fig1.png}{9.0cm}{0bp 0bp 0bp 12bp}
\bitmapfigure[max width=0.58\linewidth]{img/2010/anhui_science/q18_fig1.png}
\end{problem}

\begin{solution}
不妨取 \(EF=1\)，则 \(AB=2\). 建立空间直角坐标系，使
\(A(0,0,0)\)，\(B(2,0,0)\)，\(C(2,2,0)\)，\(D(0,2,0)\).
由 \(EF\parallel AB\)，可设 \(E(r,s,t)\)，\(F(r+1,s,t)\). 条件 \(EF\perp FB\) 给出
\[
(F-E)\cdot(B-F)=(1,0,0)\cdot(1-r,-s,-t)=0
\]
所以 \(r=1\). 又 \(BF=FC\) 给出
\[
s^2+t^2=(2-s)^2+t^2
\]
从而 \(s=1\). 再由 \(\angle BFC=90^\circ\)，
\[
(B-F)\cdot(C-F)=(0,-1,-t)\cdot(0,1,-t)=-1+t^2=0
\]
因为 \(F\) 在正方形所在平面的上方，\(t=1\). 因此
\(E(1,1,1)\)，\(F(2,1,1)\)，\(H(2,1,0)\).
平面 \(EDB\) 的一个法向量为
\[
\overrightarrow{BD}\times\overrightarrow{BE}
=(-2,2,0)\times(-1,1,1)=(2,2,0)
\]
而 \(\overrightarrow{FH}=(0,0,-1)\) 与该法向量的数量积为 \(0\)，且直线 \(FH\) 上各点满足 \(x+y=3\)，平面 \(EDB\) 的方程为 \(x+y=2\)，二者不相交. 所以 \(FH\parallel\) 平面 \(EDB\).

又
\[
\overrightarrow{AC}=(2,2,0)
\]
与平面 \(EDB\) 的法向量平行，所以 \(AC\perp\) 平面 \(EDB\).

平面 \(BDE\) 的一个法向量可取 \(\bs n_1=(1,1,0)\)，平面 \(CDE\) 的一个法向量可取
\[
\bs n_2=\overrightarrow{CD}\times\overrightarrow{CE}
=(-2,0,0)\times(-1,-1,1)=(0,2,2)
\]
故二面角 \(B-DE-C\) 的平面角 \(\varphi\) 满足
\[
\cos\varphi=\frac{\bs n_1\cdot(0,1,1)}
{|\bs n_1|\sqrt{1^2+1^2}}
=\frac{1}{2}
\]
取二面角的锐角，得 \(\varphi=60^\circ\).
\end{solution}

\begin{problem}
如图，已知椭圆 \(E\) 经过点 \(A(2,3)\)，对称轴为坐标轴，焦点 \(F_{1}\)，\(F_{2}\) 在 \(x\) 轴上，离心率 \(e = \frac{1}{2}\).

\begin{enumerate}
\item 求椭圆 \(E\) 的方程；
\item 求 \(\angle F_{1}AF_{2}\) 的角平分线所在直线 \(l\) 的方程；
\item 在椭圆 \(E\) 上是否存在关于直线 \(l\) 对称的相异两点？ 若存在，请找出；若不存在，说明理由.
\end{enumerate}

\bitmapfigure[max width=0.44\linewidth]{img/2010/anhui_science/q19_fig1.png}
\end{problem}

\begin{solution}
\begin{enumerate}
\item 设椭圆方程为
\(\frac{x^2}{a^2}+\frac{y^2}{b^2}=1\)，\(a>b>0\)
焦距的一半为 \(c\). 由离心率 \(e=\frac{c}{a}=\frac{1}{2}\)，得 \(c=\frac{a}{2}\)，从而
\[
b^2=a^2-c^2=\frac{3a^2}{4}
\]
将 \(A(2,3)\) 代入，得
\(\frac{4}{a^2}+\frac{9}{3a^2/4}=1 \Longrightarrow \frac{16}{a^2}=1 \Longrightarrow a^2=16\)，\(b^2=12\).
所以椭圆方程为
\[
\frac{x^2}{16}+\frac{y^2}{12}=1
\]
\item 由 \(c=2\)，得 \(F_1(-2,0)\)，\(F_2(2,0)\). 直线 \(AF_1\) 和 \(AF_2\) 的方程分别为
\(3x-4y+6=0\)，\(x-2=0\).
角平分线上的点到这两条直线的距离相等. 取经过椭圆内部的内角平分线，满足
\[
\frac{3x-4y+6}{5}=-(x-2)
\]
化简得
\[
l:\ 2x-y-1=0
\]
\item 假设存在椭圆上关于 \(l\) 对称的相异两点 \(P\)，\(Q\). 设它们的中点为 \(M(u,v)\). 因为 \(l\) 是线段 \(PQ\) 的垂直平分线，所以 \(M\in l\)，且可设
\(P=M+t(2,-1)\)，\(Q=M-t(2,-1)\)，\(t\neq0\).
将 \(P\)，\(Q\) 代入椭圆方程并相减，得
\[
\frac{u}{2}t-\frac{v}{3}t=0
\]
所以 \(3u-2v=0\). 又 \(M\in l\)，即 \(v=2u-1\)，联立得 \(M=(2,3)=A\).
\[
\frac{(u+2t)^2}{16}+\frac{(v-t)^2}{12}
=\frac{u^2}{16}+\frac{v^2}{12}
+t\left(\frac{u}{4}-\frac{v}{6}\right)
+t^2\left(\frac{1}{4}+\frac{1}{12}\right)
\]
在 \(M=(2,3)\) 时，中间一次项为 \(0\)，且 \(M\) 在椭圆上，所以右端等于
\[
1+\frac{t^2}{3}>1
\]
这与 \(P\) 在椭圆上要求左端等于 \(1\) 矛盾. 因此不存在关于 \(l\) 对称的相异两点.
\end{enumerate}
\end{solution}

\begin{problem}
设数列 \(a_{1}\)，\(a_{2}\)，\(\cdots\)，\(a_{n}\)，\(\cdots\) 中的每一项都不为 \(0\). 证明： \(\{a_{n}\}\) 为等差数列的充分必要条件是：对任何 \(n \in \N^{+}\)，都有 \(\frac{1}{a_{1} a_{2}} + \frac{1}{a_{2} a_{3}} + \cdots + \frac{1}{a_{n} a_{n+1}} = \frac{n}{a_{1} a_{n+1}}\).
\end{problem}

\begin{solution}
先证充分性. 若 \(\{a_n\}\) 为等差数列，设公差为 \(d\). 当 \(d\neq0\) 时，
\[
\frac{1}{a_k a_{k+1}}
=\frac{1}{d}\left(\frac{1}{a_k}-\frac{1}{a_{k+1}}\right)
\]
于是
\[
\sum_{k=1}^{n}\frac{1}{a_k a_{k+1}}
=\frac{1}{d}\left(\frac{1}{a_1}-\frac{1}{a_{n+1}}\right)
=\frac{a_{n+1}-a_1}{d a_1a_{n+1}}
=\frac{n}{a_1a_{n+1}}
\]
当 \(d=0\) 时，所有项都等于非零常数 \(a_1\)，等式两边均为 \(\frac{n}{a_1^2}\)，结论同样成立.

再证必要性. 设题中等式对任意 \(n\in\N^+\) 成立. 对 \(n\geqslant2\)，用第 \(n\) 个等式减去第 \(n-1\) 个等式，得
\[
\frac{1}{a_na_{n+1}}
=\frac{n}{a_1a_{n+1}}-\frac{n-1}{a_1a_n}
\]
由于各项均不为 \(0\)，两边同乘 \(a_1a_na_{n+1}\)，得
\(a_1=na_n-(n-1)a_{n+1}\)，\(a_{n+1}=\frac{na_n-a_1}{n-1}\).
令 \(d=a_2-a_1\). 当 \(n=2\) 时，递推式给出 \(a_3=2a_2-a_1=a_1+2d\). 若 \(a_n=a_1+(n-1)d\)，则
\[
a_{n+1}
=\frac{n[a_1+(n-1)d]-a_1}{n-1}
=a_1+nd
\]
由数学归纳法，\(a_n=a_1+(n-1)d\) 对所有 \(n\in\N^+\) 成立，故 \(\{a_n\}\) 为等差数列.
\end{solution}

\begin{problem}
品酒师需定期接受酒味鉴别功能测试，一般通常采用的测试方法如下：拿出 \(n\) 瓶外观相同但品质不同的酒让其品尝，要求其按品质优劣为它们排序；经过一段时间，等其记忆淡忘之后，再让其品尝这 \(n\) 瓶酒，并重新按品质优劣为它们排序，这称为一轮测试. 根据一轮测试中的两次排序的偏离程度的高低为其评分.

现设 \(n = 4\)，分别以 \(a_{1}\)，\(a_{2}\)，\(a_{3}\)，\(a_{4}\) 表示第一次排序时被排为 \(1\)，\(2\)，\(3\)，\(4\) 的四种酒在第二次排序时的序号，并令 \(X = |1 - a_{1}| + |2 - a_{2}| + |3 - a_{3}| + |4 - a_{4}|\)，则 \(X\) 是对两次排序的偏离程度的一种描述.

\begin{enumerate}
\item 写出 \(X\) 的可能值集合；
\item 假设 \(a_{1}\)，\(a_{2}\)，\(a_{3}\)，\(a_{4}\) 等可能地为 \(1\)，\(2\)，\(3\)，\(4\) 的各种排列，求 \(X\) 的分布列；
\item 某品酒师在相继进行的三轮测试中，都有 \(X \leqslant 2\)，
\begin{enumerate}
\item 试按（2）中的结果，计算出现这种现象的概率 （假定各轮测试相互独立）；
\item 你认为该品酒师的酒味鉴别功能如何？ 说明理由.
\end{enumerate}
\end{enumerate}
\end{problem}

\begin{solution}
\begin{enumerate}
\item 对任意排列，\(|i-a_i|\equiv i-a_i\pmod 2\). 因为 \(\sum_{i=1}^{4}i=\sum_{i=1}^{4}a_i=10\)，所以 \(X\) 必为偶数. 又 \(0\leqslant X\leqslant8\)，并且下列排列分别给出 \(0\)，\(2\)，\(4\)，\(6\)，\(8\)，所以
\[
X\text{ 的可能值集合为 }\{0,2,4,6,8\}
\]
\item 逐一计算 \(24\) 个排列的偏离值，可按 \(X\) 分组如下：
当 \(X=0\) 时，排列为 \((1,2,3,4)\).
当 \(X=2\) 时，排列为 \((1,2,4,3)\)、\((1,3,2,4)\)、\((2,1,3,4)\).
当 \(X=4\) 时，排列为 \((1,3,4,2)\)、\((1,4,2,3)\)、\((1,4,3,2)\)、\((2,1,4,3)\)、\((2,3,1,4)\)、\((3,1,2,4)\)、\((3,2,1,4)\).
当 \(X=6\) 时，排列为 \((2,3,4,1)\)、\((2,4,1,3)\)、\((2,4,3,1)\)、\((3,1,4,2)\)、\((3,2,4,1)\)、\((4,1,2,3)\)、\((4,1,3,2)\)、\((4,2,1,3)\)、\((4,2,3,1)\).
当 \(X=8\) 时，排列为 \((3,4,1,2)\)、\((3,4,2,1)\)、\((4,3,1,2)\)、\((4,3,2,1)\).
因此
\(P(X=0)=\frac{1}{24}\)，\(P(X=2)=\frac{3}{24}\)，\(P(X=4)=\frac{7}{24}\)，\(P(X=6)=\frac{9}{24}\)，\(P(X=8)=\frac{4}{24}\).
\item
\begin{enumerate}
\item 一轮测试中 \(X\leqslant2\) 的概率为
\[
P(X\leqslant2)=P(X=0)+P(X=2)=\frac{1+3}{24}=\frac{1}{6}
\]
三轮相互独立，所以三轮均满足 \(X\leqslant2\) 的概率为
\[
\left(\frac{1}{6}\right)^3=\frac{1}{216}
\]
\item 若第二次排序完全随机，这种三轮结果出现的概率仅为 \(\frac{1}{216}\approx0.0046\)，概率很小，因此可认为该品酒师具有较好的酒味鉴别功能；这说明测试结果支持其排序能力较强，但不能仅凭三轮测试作绝对结论.
\end{enumerate}
\end{enumerate}
\end{solution}
